\section{Introduction}
\label{sec:intro}

In-context learning (ICL) recasts medical image segmentation as a
few-shot matching problem: given a target scan and a small set of
(image, mask) context pairs illustrating the structure of interest, the
model predicts a mask for the target without any task-specific
fine-tuning. This removes the need to retrain or fine-tune a model for
every new anatomical class, which is especially valuable for structures
that are rare, newly defined, or only annotated by a handful of clinical
sites.

Most ICL segmentation models were designed and validated in
2D~\cite{butoiUniverSegUniversalMedical2023,gaoShowSegmentUniversal2025,huMedverseUniversalModel2025a}.
Moving to native 3D volumes is attractive -- CT and MRI are inherently
volumetric, and slice-wise processing discards through-plane context
that is often necessary to resolve thin or elongated structures -- but
it sharpens an already difficult resolution--compute trade-off. Dense
cross-attention between a target volume and $K$ context volumes scales
with the cube of the side length, so a design that is affordable at
$64^3$ becomes impractical at $256^3$. Existing 3D-capable approaches
each make a different concession: Medverse~\cite{huMedverseUniversalModel2025a}
processes the whole volume through a sliding window at every resolution
step of its autoregressive cascade, re-visiting the entire volume even
once coarse levels have already localized the target; Iris~\cite{gaoShowSegmentUniversal2025}
fuses context information into the target through a single-shot
pixel-shuffle/unshuffle operation with no iterative refinement across
scales. Our own concurrent 2D work on hierarchical patch selection~\cite{camaretPatchICL2026}
showed that supervising
which regions to process, rather than relying on attention alone to
discover them, yields a favorable accuracy/compute trade-off in 2D; here
we ask whether the same principle extends to full 3D volumes, and how
the image--label interaction itself should be designed once both axes
(space and image/label identity) matter jointly.

We propose \textbf{PatchICL-3D}, an in-context 3D segmentation model
with two components that address these gaps directly.

\vspace{0.3em}\noindent\textbf{Bi-axial image--label attention.}
Instead of concatenating image and mask channels before a shared
encoder (Medverse) or fusing them once via pixel shuffle-unshuffle
(Iris), we let target and context tokens attend jointly along two axes:
a spatial-patch axis (as in standard within-volume/cross-context
attention) and an image/label axis that lets label information at a
given location refine, and be refined by, the corresponding image
features. This keeps the image and label streams distinct for longer
while still allowing early, repeated exchange between them.

\vspace{0.3em}\noindent\textbf{Register-carried coarse-to-fine cascade.}
PatchICL-3D processes a volume through a small number of resolution
levels. Two design choices connect the levels: (1) a \emph{query image
prior} -- the previous level's upsampled prediction initializes the
target's label tokens at the next level, rather than entering the
network as a separate perturbed-GT branch (Medverse) or not at all
(Iris); and (2) \emph{region-restricted fine processing} -- once a
coarse level localizes the structure, the sliding window at finer levels
is restricted to the predicted zone instead of re-scanning the full
volume. To let the cascade be trained as a single model rather than a
chain of independently-supervised stages, we introduce a small number of
learned \emph{register tokens} that carry hidden state from one level's
processing into the next, giving the optimizer a gradient path that
crosses resolution levels.

\vspace{0.3em}\noindent\textbf{Contributions.} (1) A bi-axial
image--label attention mechanism for in-context 3D segmentation that
attends jointly over space and image/label identity, contrasted directly
against early-concatenation and one-shot fusion baselines. (2) A
coarse-to-fine cascade in which the query prior and the region
restriction at fine levels are ablated independently, quantifying their
individual contribution to the accuracy/compute trade-off. (3)
Register-token cross-level training that lets gradients flow across the
cascade, evaluated against per-level-only supervision. (4) An evaluation
on TotalSegmentator~\cite{wasserthalTotalSegmentatorRobustSegmentation2023}
at fixed spacing against Medverse and Iris on accuracy (Dice, NSD) and
compute (time, FLOPs, VRAM), together with generalization experiments on
\TODO{other CT datasets, other modalities, and far-OOD tasks}.
